% c08 原书 p.91–103（PDF p.098–110）
\subsection{概形与层的形式完全化}
\label{sec:1.5}

本节的主要内容只在本书的最后一节用到。

设 $(M_n,\phi_{mn})_{n\in\bN}$ 是集合的逆向系统。若对每个 $n$，存在自然数 $n_0$，使得对每个 $m\ge n+n_0$ 都有
\[
\Image\bigl(M_{n+n_0}\xrightarrow{\ \phi_{n+n_0,n}\ }M_n\bigr)
=\Image\bigl(M_m\xrightarrow{\ \phi_{mn}\ }M_n\bigr),
\]
则称它满足 \emph{Mittag-Leffler 条件}。
\begin{tnote}
直观地说，Mittag-Leffler 条件要求：对每个 $n$，当 $m$ 不断增大时，$M_m$ 在 $M_n$ 中的像很快就稳定下来。这个条件排除了「像无限地逐次变小」的情形，因而保证下面 \ref{it:1.5.1} 中的逆向极限序列仍保持正合。[依据：原书 p.92 命题 1.5.1 及其证明中对本条件的用法]
\end{tnote}

设 $(M_n,\phi_{mn})_{n\in\bN}$ 满足 Mittag-Leffler 条件，且每个 $M_n$ 非空。令 $M'_n=\bigcap_{m\ge n}\Image\phi_{mn}$。则 $(M'_n)_{n\in\bN}$ 是非空集合的逆向系统，且
\[
\plim_n M'_n\cong\plim_n M_n.
\]
此外，每个 $M'_m\xrightarrow{\phi_{mn}}M'_n$ 都是满射。由此推出 $\plim_n M'_n\ne\emptyset$。

注意，若每个 $\phi_{mn}$ 都是满射，则 $(M_n,\phi_{mn})_{n\in\bN}$ 满足 Mittag-Leffler 条件。若 $(M_n,\phi_{mn})_{n\in\bN}$ 是 $A$-模的逆向系统，且每个 $M_n$ 有有限长度，则 $(M_n)_{n\in\bN}$ 满足 Mittag-Leffler 条件。事实上，对每个 $n$，递减链
\[
\Image\phi_{nn}\supset\Image\phi_{n+1,n}\supset\Image\phi_{n+2,n}\supset\cdots
\]
终将稳定。

\begin{proposition}\label{it:1.5.1}
设 $(M'_n,\phi'_{mn})_{n\in\bN}$、$(M_n,\phi_{mn})_{n\in\bN}$ 与 $(M''_n,\phi''_{mn})_{n\in\bN}$ 是交换群的逆向系统，$(M'_n\xrightarrow{u_n}M_n)_{n\in\bN}$ 与 $(M_n\xrightarrow{v_n}M''_n)_{n\in\bN}$ 是逆向系统的态射，使得对每个 $n$，序列
\[
0\to M'_n\xrightarrow{u_n}M_n\xrightarrow{v_n}M''_n\to 0
\]
正合。若 $(M'_n)_{n\in\bN}$ 满足 Mittag-Leffler 条件，则序列
\[
0\to\plim_n M'_n\to\plim_n M_n\to\plim_n M''_n\to 0
\]
正合。
\end{proposition}

\begin{proof}
只有 $\plim_n M_n\to\plim_n M''_n$ 的满射性不是平凡的。给定 $(x''_n)\in\plim_n M''_n$，令 $N_n=v_n^{-1}(x''_n)$。则 $(N_n,\phi_{mn})_{n\in\bN}$ 是集合的逆向系统。由于每个 $v_n$ 都是满射，每个 $N_n$ 都非空。我们断言 $(N_n)_{n\in\bN}$ 满足 Mittag-Leffler 条件。于是由本节开头的讨论，$\plim_n N_n\ne\emptyset$。$\plim_n N_n$ 中每个元素都属于 $\plim_n M_n$，并被打到 $(x''_n)$。这就证明了 $\plim_n M_n\to\plim_n M''_n$ 是满射。下面证明上述断言。给定 $n$，由于 $(M'_n)_{n\in\bN}$ 满足 Mittag-Leffler 条件，可找到自然数 $n_0$，使对每个 $m\ge n+n_0$ 都有
\[
\Image(M'_{n+n_0}\xrightarrow{\phi'_{n+n_0,n}}M'_n)=\Image(M'_m\xrightarrow{\phi'_{mn}}M'_n).
\]
固定 $m\ge n+n_0$ 与元素 $x_m\in N_m$。对每个 $n\le k\le m$，令 $x_k=\phi_{mk}(x_m)$。则映射族
\[
M'_k\to N_k,\quad x'_k\mapsto u_k(x'_k)+x_k
\]
定义了集合的逆向系统 $(M'_k)_{n\le k\le m}$ 与 $(N_k)_{n\le k\le m}$ 之间的同构。所以有
\[
\Image(N_{n+n_0}\xrightarrow{\phi_{n+n_0,n}}N_n)=\Image(N_m\xrightarrow{\phi_{mn}}N_n).
\]
这对每个 $m\ge n+n_0$ 都成立。所以 $(N_n)_{n\in\bN}$ 满足 Mittag-Leffler 条件。
\end{proof}
\begin{tnote}
一般说来，逆向极限只是左正合的，并不保持满射。\ref{it:1.5.1} 说明：只要核系统 $(M'_n)$ 满足 Mittag-Leffler 条件，逆向极限就仍然保持短正合序列。下面 \ref{it:1.5.3}、\ref{it:1.5.4} 与 \ref{it:1.5.13} 都要用到这一点。[依据：原书 p.92–95，1.5.1 之后各命题的证明]
\end{tnote}

设 $A$ 是环，$I$ 是 $A$ 的理想，$M$ 是 $A$-模。$M$ 的递减滤链
\[
M=M_0\supset M_1\supset M_2\supset\cdots
\]
若对每个 $n$ 都有 $IM_n\subset M_{n+1}$，则称它为 \emph{$I$-滤链}\idx{ijinlvlian}{$I$-滤链}。由此推出 $I^nM\subset M_n$ 对每个 $n$ 成立。若它是 $I$-滤链，且存在自然数 $n_0$，使对每个 $n\ge n_0$ 都有 $IM_n=M_{n+1}$，则称该滤链为 \emph{$I$-稳定}\idx{ijinwending}{$I$-稳定}的。由此推出，对每个 $n\ge n_0$，
\[
I^nM\subset M_n=I^{n-n_0}M_{n_0}\subset I^{n-n_0}M,
\]
从而
\[
\plim_n M/I^nM\cong\plim_n M/M_n.
\]

\begin{proposition}\label{it:1.5.2}
设 $A$ 是环，$I$ 是 $A$ 的理想，$M$ 是 $A$-模，
\[
M=M_0\supset M_1\supset M_2\supset\cdots
\]
是 $M$ 的 $I$-滤链，$S=\bigoplus_{n=0}^{\infty}I^n$ 是分次环。设每个 $M_n$ 都是有限生成 $A$-模。则此滤链是 $I$-稳定的，当且仅当分次 $S$-模 $\bigoplus_{n=0}^{\infty}M_n$ 是有限生成的。
\end{proposition}

\begin{proof}
由于该滤链是 $I$-滤链，我们有 $I^nM_m\subset M_{m+n}$。这在 $\bigoplus_{n=0}^{\infty}M_n$ 上定义了分次 $S$-模结构。设滤链是 $I$-稳定的。取自然数 $n_0$，使 $M_n=I^{n-n_0}M_{n_0}$ 对一切 $n\ge n_0$ 成立。则作为 $S$-模，$\bigoplus_{n=0}^{\infty}M_n$ 可由 $\bigoplus_{n=0}^{n_0}M_n$ 生成。由于每个 $M_n$ 是有限生成 $S_0$-模，$\bigoplus_{n=0}^{\infty}M_n$ 是有限生成 $S$-模。

反之，设 $\bigoplus_{n=0}^{\infty}M_n$ 是有限生成 $S$-模。设 $x_i$（$i=1,\dots,k$）是一族生成 $\bigoplus_{n=0}^{\infty}M_n$ 的齐次元素，令 $n_i=\deg x_i$，$n_0=\max(n_1,\dots,n_k)$。对每个 $n\ge n_0$，每个 $x\in M_n$ 都可写成 $r_1x_1+\cdots+r_kx_k$，其中 $r_i\in S_{n-n_i}=I^{n-n_i}$。我们可把每个 $r_i$ 写成 $\sum_j a_{ij}r_{ij}$，其中 $a_{ij}\in I^{n-n_0}$，$r_{ij}\in I^{n_0-n_i}$。于是
\[
x=\sum_{i,j}a_{ij}(r_{ij}x_i)\in I^{n-n_0}M_{n_0}.
\]
所以对每个 $n\ge n_0$ 都有 $M_n=I^{n-n_0}M_{n_0}$，从而该滤链是 $I$-稳定的。
\end{proof}

\begin{corollary}[Artin-Rees]\label{it:1.5.3}
设 $A$ 是诺特环，$I$ 是 $A$ 的理想，$M$ 是有限生成 $A$-模，$N$ 是 $M$ 的子模。则 $N$ 的滤链
\[
N\supset N\cap IM\supset N\cap I^2M\supset\cdots
\]
是 $I$-稳定的，从而
\[
\plim_n N/(N\cap I^nM)\cong\plim_n N/I^nN.
\]
\end{corollary}

\begin{proof}
该滤链显然是 $I$ 的 $I$-滤链。
\begin{tnote}
原书 p.94 此处作「an $I$-filtration of $I$」，应为「$N$ 的 $I$-滤链」；所论滤链正是 $N\supset N\cap IM\supset N\cap I^2M\supset\cdots$。[依据：原书 p.94 页图像，参看同页 Corollary 1.5.3 的陈述]
\end{tnote}
分次环 $S=\bigoplus_{n=0}^{\infty}I^n$ 是诺特的。事实上，若 $r_1,\dots,r_k$ 是理想 $I$ 的一族生成元，则 $A$-代数同态
\[
A[x_1,\dots,x_k]\to S,\quad x_i\mapsto r_i\in S_1
\]
是满射。我们的断言于是由 $A[x_1,\dots,x_k]$ 是诺特的推出。由于
\[
M\supset IM\supset I^2M\supset\cdots
\]
是 $M$ 的 $I$-稳定滤链，分次 $S$-模 $\bigoplus_{n=0}^{\infty}I^nM$ 是有限生成的。所以它的子模 $\bigoplus_{n=0}^{\infty}(N\cap I^nM)$ 也是有限生成的。于是由 \ref{it:1.5.2}，$N$ 的 $I$-滤链
\[
N\supset N\cap IM\supset N\cap I^2M\supset\cdots
\]
是 $I$-稳定的。
\end{proof}

设 $A$ 是环，$I$ 是 $A$ 的理想。环
\[
\widehat A=\plim_n A/I^n
\]
称为 $A$ 的 \emph{$I$-进完全化}\idx{ijinwanquanhua}{$I$-进完全化}。有典范同态
\[
A\to\widehat A.
\]
若此同态是同构，则称 $A$ 是 \emph{$I$-进完备}\idx{ijinwanbei}{$I$-进完备}的。对任意 $A$-模 $M$，$A$-模
\[
\widehat M=\plim_n M/I^nM
\]
称为 $M$ 的 \emph{$I$-进完全化}。有典范同态
\[
M\to\widehat M,
\]
它与模结构及环同态 $A\to\widehat A$ 相容。若 $M\to\widehat M$ 是同构，则称 $M$ 是 \emph{$I$-进完备}的。

\begin{proposition}\label{it:1.5.4}
设 $A$ 是诺特环，$I$ 是 $A$ 的理想，
\[
0\to M'\to M\to M''\to 0
\]
是有限生成 $A$-模的正合序列。则序列
\[
0\to\widehat{M'}\to\widehat M\to\widehat{M''}\to 0
\]
正合。
\end{proposition}

\begin{proof}
对每个 $n$，有正合序列
\[
0\to M'/(M'\cap I^nM)\to M/I^nM\to M''/I^nM''\to 0.
\]
取逆向极限并应用 \ref{it:1.5.1}，即知
\[
0\to\plim_n M'/(M'\cap I^nM)\to\plim_n M/I^nM\to\plim_n M''/I^nM''\to 0
\]
正合。由 \ref{it:1.5.3}，此序列正是
\[
0\to\widehat{M'}\to\widehat M\to\widehat{M''}\to 0.
\]
\end{proof}

\begin{proposition}\label{it:1.5.5}
设 $A$ 是诺特环，$I$ 是 $A$ 的理想，$M$ 是有限生成 $A$-模，$\widehat A$ 与 $\widehat M$ 分别是 $A$ 与 $M$ 的 $I$-进完全化。则有典范同构
\[
\widehat A\tensor_A M\cong\widehat M.
\]
特别地，若 $A$ 是 $I$-进完备的，则任意有限生成 $A$-模 $M$ 也是 $I$-进完备的。
\end{proposition}

\begin{proof}
同态 $M\to\widehat M$ 诱导同态
\[
\widehat A\tensor_A M\to\widehat M.
\]
由于 $M$ 是有限生成 $A$-模且 $A$ 是诺特的，可找到正合序列
\[
A^m\to A^n\to M\to 0.
\]
考虑交换图
\begin{center}
\begin{tikzcd}
\widehat A\tensor_A A^m \arrow[r] \arrow[d] & \widehat A\tensor_A A^n \arrow[r] \arrow[d] & \widehat A\tensor_A M \arrow[r] \arrow[d] & 0\\
(\widehat A)^m \arrow[r] & (\widehat A)^n \arrow[r] & \widehat M \arrow[r] & 0.
\end{tikzcd}
\end{center}
第一行正合，因为函子 $\widehat A\tensor_A -$ 是右正合的。第二行由 \ref{it:1.5.4} 正合。前两个竖直箭头是同构。所以最后一个竖直箭头也是同构。
\end{proof}

\begin{corollary}\label{it:1.5.6}
设 $A$ 是诺特环，$I$ 是 $A$ 的理想，$\widehat A$ 是 $A$ 的 $I$-进完全化。
\begin{enumerate}[label=(\roman*),leftmargin=2.8em]
  \item 对 $A$ 的每个理想 $J$，$J$ 的 $I$-进完全化 $\widehat J$ 是 $\widehat A$ 的理想。
  \item 对每个 $n$，有 $A/I^n\cong\widehat A/\widehat I^n$。因此 $\widehat I^n$ 是典范同态 $\widehat A\to A/I^n$ 的核。
  \item 对每对 $0\le m\le n$，有 $\widehat I^n\subset\widehat I^m$ 且 $I^m/I^n\cong\widehat I^m/\widehat I^n$。
  \item 对 $A$ 的每个理想 $J$ 与每个 $n$，有
  \[
  \widehat{J^n}\cong\widehat A\tensor_A J^n\cong J^n\widehat A=(J\widehat A)^n\cong\widehat J^{\,n}.
  \]
\end{enumerate}
\end{corollary}

\begin{proof}
由 \ref{it:1.5.4}，序列
\[
0\to\widehat J\to\widehat A
\]
正合。所以 $\widehat J$ 是 $\widehat A$ 的理想。对每对 $0\le m\le n$，有正合序列
\[
0\to I^n\to I^m\to I^m/I^n\to 0.
\]
由 \ref{it:1.5.4}，序列
\[
0\to\widehat{I^n}\to\widehat{I^m}\to\widehat{I^m/I^n}\to 0
\]
正合。由 $I$-进完全化的定义可证
\[
\widehat{I^m/I^n}\cong I^m/I^n.
\]
所以 $\widehat{I^n}\subset\widehat{I^m}$ 且 $\widehat{I^m}/\widehat{I^n}\cong I^m/I^n$。取 $m=0$，得
\[
\widehat A/\widehat{I^n}\cong A/I^n.
\]
由 \ref{it:1.5.4} 与 \ref{it:1.5.5}，序列
\[
0\to\widehat A\tensor_A J^n\to\widehat A\tensor_A A\to\widehat A\tensor_A A/J^n\to 0
\]
正合。特别地，典范同态 $\widehat A\tensor_A J^n\to\widehat A$ 是单射。其像为 $J^n\widehat A$。所以
\[
\widehat A\tensor_A J^n\cong J^n\widehat A.
\]
(iv) 中的其余断言于是显然。
\end{proof}

\begin{proposition}\label{it:1.5.7}
设 $A$ 是诺特环，$I$ 是 $A$ 的理想，$\widehat A$ 是 $A$ 的 $I$-进完全化。则 $\widehat A$ 是 $\widehat I$-进完备的，且 $\widehat I\subset\operatorname{rad}(\widehat A)$。
\end{proposition}

回忆：环 $A$ 的 \emph{Jacobson 根}\idx{jacobsongeng}{Jacobson 根} $\operatorname{rad}(A)$ 定义为 $A$ 的所有极大理想的交。

\begin{proof}
我们有
\[
\widehat A=\plim_n A/I^n\cong\plim_n\widehat A/\widehat I^n.
\]
所以 $\widehat A$ 是 $\widehat I$-进完备的。对每个 $x\in\widehat I$ 与每个 $r\in\widehat A$，考虑级数
\[
1+rx+(rx)^2+\cdots.
\]
注意它在每个 $\widehat A/\widehat I^n$ 中的像有意义，因为它是有限和
\[
1+rx+(rx)^2+\cdots+(rx)^{n-1}.
\]
所以 $1+rx+(rx)^2+\cdots$ 定义了 $\widehat A=\plim_n\widehat A/\widehat I^n$ 中的一个元素。乘积 $(1-rx)(1+rx+(rx)^2+\cdots)$ 在每个 $\widehat A/\widehat I^n$ 中的像是 $1$。所以
\[
(1-rx)(1+rx+(rx)^2+\cdots)=1,
\]
从而 $1-rx$ 是 $\widehat A$ 中的单位。设 $\fm$ 是 $\widehat A$ 的极大理想。若 $x\notin\fm$，则 $(x)+\fm=\widehat A$，从而 $rx+y=1$，其中 $r\in\widehat A$，$y\in\fm$。但此时 $1-rx=y$ 不可能是单位。矛盾。所以 $x$ 属于 $\widehat A$ 的每个极大理想。
\end{proof}

\begin{proposition}\label{it:1.5.8}
设 $A$ 是诺特环，$I$ 是 $A$ 的理想，$M$ 是有限生成 $A$-模，$\widehat M$ 是 $M$ 的 $I$-进完全化。则
\[
\ker(M\to\widehat M)=\bigcap_{n=1}^{\infty}I^nM=\{x\in M\mid (1+r)x=0 \text{ 对某个 } r\in I\}.
\]
若 $I\subset\operatorname{rad}(A)$，则 $\bigcap_{n=1}^{\infty}I^nM=0$。
\end{proposition}

\begin{proof}
若 $x\in M$ 且对某个 $r\in I$ 有 $(1+r)x=0$，则
\[
x=(-r)x=(-r)^2x=\cdots\in\bigcap_{n=1}^{\infty}I^nM.
\]
令 $N=\ker(M\to\widehat M)=\bigcap_{n=1}^{\infty}I^nM$。对每个 $n$ 有 $N\cap I^nM=N$。由 \ref{it:1.5.3}，必有 $IN=N$。设 $x_1,\dots,x_k$ 是 $N$ 的一族有限生成元。则可写成
\[
x_i=\sum_j a_{ij}x_j
\]
其中 $a_{ij}\in I$。设 $P$ 是矩阵 $(a_{ij})$，$E$ 是单位矩阵。则有
\[
(E-P)\begin{pmatrix}x_1\\ \vdots\\ x_k\end{pmatrix}=0.
\]
设 $(E-P)^*$ 是 $E-P$ 的伴随矩阵。则
\[
(E-P)^*(E-P)\begin{pmatrix}x_1\\ \vdots\\ x_k\end{pmatrix}=0.
\]
所以对每个 $i$ 有 $\det(E-P)x_i=0$。由于 $P$ 的元素都在 $I$ 中，有 $\det(E-P)=1+r$，其中 $r\in I$。于是对每个 $i$ 有 $(1+r)x_i=0$，从而对每个 $x\in N=\bigcap_{n=1}^{\infty}I^nM$ 有 $(1+r)x=0$。再设 $I\subset\operatorname{rad}A$。则 $1+r$ 是 $A$ 中的单位。所以 $N=0$。
\end{proof}
\begin{tnote}
\ref{it:1.5.8} 说明，$I$-进完全化丢掉的正是那些被某个 $1+r$（$r\in I$）零化的元素。因此当 $I$ 落在 Jacobson 根中时，完全化映射 $M\to\widehat M$ 是单射；但在一般情况下它可能不是单射。[依据：原书 p.97–98 命题 1.5.8 及其证明]
\end{tnote}

\begin{proposition}\label{it:1.5.9}
设 $A$ 是诺特环，$I$ 是 $A$ 的理想，$\widehat A$ 是 $A$ 的 $I$-进完全化。则 $\widehat A$ 是诺特的。
\end{proposition}

\begin{proof}
考虑分次环
\[
G(\widehat A)=\bigoplus_{n=0}^{\infty}\widehat I^n/\widehat I^{n+1}\cong\bigoplus_{n=0}^{\infty}I^n/I^{n+1}.
\]
它是诺特的。事实上，设 $r_1,\dots,r_m$ 是理想 $I$ 的一族有限生成元，$\bar r_1,\dots,\bar r_m$ 是它们在 $I/I^2$ 中的像。则 $A/I$-代数同态
\[
A/I[x_1,\dots,x_m]\to G(\widehat A),\quad x_i\mapsto\bar r_i
\]
是满射。我们的断言于是由 $A/I[x_1,\dots,x_m]$ 是诺特的推出。

设 $J$ 是 $\widehat A$ 的理想。我们要证它是有限生成的。考虑分次 $G(\widehat A)$-模
\[
G(J)=\bigoplus_{n=0}^{\infty}J\cap\widehat I^n/J\cap\widehat I^{n+1}.
\]
它是 $G(\widehat A)$ 的分次子模。由于 $G(\widehat A)$ 是诺特的，$G(J)$ 是有限生成的。设 $a_i\in J\cap\widehat I^{n_i}$（$i=1,\dots,k$），使得它们在 $G(J)$ 中的像构成一族生成元。设 $\widehat A_i$ 是 $\widehat A$ 自身配上 $\widehat I$-稳定滤链
\[
(\widehat A_i)_n=\begin{cases}
\widehat A, & \text{若 } 0\le n\le n_i,\\
\widehat I^{n-n_i}, & \text{若 } n\ge n_i.
\end{cases}
\]
考虑 $\widehat A$-模同态
\[
\phi\colon\bigoplus_{i=1}^{k}\widehat A_i\to J,\quad e_i\mapsto a_i,
\]
其中对每个 $i$，$e_i$ 是 $\bigoplus_{i=1}^{k}\widehat A_i$ 中第 $i$ 个分量为 $1$、其余分量为 $0$ 的元素。则有
\[
\phi\Bigl(\bigl(\bigoplus_{i=1}^{k}\widehat A_i\bigr)_n\Bigr)\subset J\cap\widehat I^n.
\]
于是 $\phi$ 诱导分次 $G(\widehat A)$-模同态
\[
G(\phi)\colon G\Bigl(\bigoplus_{i=1}^{k}\widehat A_i\Bigr)
=\bigoplus_{n=0}^{\infty}\bigl(\bigoplus_{i=1}^{k}\widehat A_i\bigr)_n\Big/\bigl(\bigoplus_{i=1}^{k}\widehat A_i\bigr)_{n+1}
\to G(J)=\bigoplus_{n=0}^{\infty}J\cap\widehat I^n/J\cap\widehat I^{n+1}.
\]
由 $\phi$ 的定义，$G(\phi)$ 是满射。对每个 $y\in J$，$y$ 在 $J/J\cap\widehat I$ 中的像落在 $G(\phi)$ 的像中。所以存在 $x_0\in(\bigoplus_{i=1}^{k}\widehat A_i)_0$，使
\[
\phi(x_0)\equiv y\pmod{J\cap\widehat I}.
\]
$y-\phi(x_0)\in J\cap\widehat I$ 在 $J\cap\widehat I/J\cap\widehat I^2$ 中的像落在 $G(\phi)$ 的像中。所以存在 $x_1\in(\bigoplus_{i=1}^{k}\widehat A_i)_1$，使
\[
\phi(x_1)\equiv y-\phi(x_0)\pmod{J\cap\widehat I^2}.
\]
如此下去，对每个自然数 $n$，可找到 $x_n\in(\bigoplus_{i=1}^{k}\widehat A_i)_n$，使
\[
\phi(x_1+\cdots+x_n)\equiv y\pmod{J\cap\widehat I^{n+1}}.
\]
级数 $x=x_0+x_1+\cdots$ 定义了
\[
\bigoplus_{i=1}^{k}\widehat A_i=\plim_n\bigl(\bigoplus_{i=1}^{k}\widehat A_i\bigr)\Big/\bigl(\bigoplus_{i=1}^{k}\widehat A_i\bigr)_n
\]
中的一个元素，因为它在每个 $\bigl(\bigoplus_{i=1}^{k}\widehat A_i\bigr)/\bigl(\bigoplus_{i=1}^{k}\widehat A_i\bigr)_n$ 中的像都是有限和。级数 $x_{n+1}+x_{n+2}+\cdots$ 定义了 $\bigl(\bigoplus_{i=1}^{k}\widehat A_i\bigr)_{n+1}$ 中的一个元素，因为它在每个 $\bigl(\bigoplus_{i=1}^{k}\widehat A_i\bigr)_{n+1}/\bigl(\bigoplus_{i=1}^{k}\widehat A_i\bigr)_m$ 中的像对每个 $m\ge n+1$ 都是有限和。由于
\[
\phi(x_1+\cdots+x_n)-y,\quad \phi(x_{n+1}+x_{n+2}+\cdots)\in\widehat I^{n+1},
\]
我们有
\[
\phi(x)-y=(\phi(x_1+\cdots+x_n)-y)+\phi(x_{n+1}+x_{n+2}+\cdots)\in\bigcap_{n=1}^{\infty}\widehat I^{n+1}.
\]
但 $\bigcap_{n=1}^{\infty}\widehat I^{n+1}=0$，因为 $\widehat A\cong\plim_n\widehat A/\widehat I^n$。于是 $\phi(x)=y$，从而 $\phi$ 是满射。这就证明了 $J$ 是有限生成的。因此 $\widehat A$ 是诺特的。
\end{proof}

\begin{proposition}\label{it:1.5.10}
设 $(A_n,\phi_{mn})_{n\in\bN}$ 是环的逆向系统。对每个 $n$，设 $I_n$ 是同态 $\phi_{n1}\colon A_n\to A_1$ 的核。设以下条件成立：
\begin{enumerate}[label=(\alph*),leftmargin=2.8em]
  \item 对每对 $m\ge n$，$\phi_{mn}\colon A_m\to A_n$ 是满射，其核为 $I_m^n$。特别地，$\phi_{mm}=\id_{A_m}$ 的核为 $I_m^m=0$。
  \item $I_2/I_2^2=I_2$ 作为 $A_2/I_2\cong A_1$-模是有限生成的。
\end{enumerate}
设 $A=\plim_n A_n$。对每个 $n$，设 $\phi_n\colon A\to A_n$ 是典范同态，$I^{(n)}$ 是它的核。我们有
\begin{enumerate}[label=(\roman*),leftmargin=2.8em]
  \item 对每个 $n$，$I^{(n)}$ 是有限生成的，且 $I^{(n)}=(I^{(1)})^n$。
  \item $A\cong\plim_n A/(I^{(1)})^n$。
  \item $I^{(1)}/(I^{(1)})^2\cong I_2/I_2^2$。
\end{enumerate}
\end{proposition}

\begin{proof}
注意 (i) 蕴含 (ii) 与 (iii)。事实上，由于每个 $\phi_{mn}$ 都是满射，每个 $\phi_n\colon A\to A_n$ 也是满射。它诱导同构
\[
A/I^{(n)}\cong A_n.
\]
我们有
\[
A=\plim_n A_n\cong\plim_n A/I^{(n)}\cong\plim_n A/(I^{(1)})^n.
\]
这证明了 (ii)。有交换图
\begin{center}
\begin{tikzcd}
A/(I^{(1)})^2 \arrow[r,"\cong"] \arrow[d] & A_2 \arrow[d]\\
A/I^{(1)} \arrow[r,"\cong"'] & A_1.
\end{tikzcd}
\end{center}
第一个竖直箭头的核是 $I^{(1)}/(I^{(1)})^2$，第二个竖直箭头的核是 $I_2=I_2/I_2^2$。所以 $I^{(1)}/(I^{(1)})^2\cong I_2/I_2^2$。这证明了 (iii)。

下证 (i)。设 $a_i=(\dots,a_{i3},a_{i2},0)$（$i=1,\dots,k$）是 $A=\plim_n A_n$ 中一族有限个元素，使得 $a_{i2}$ 作为 $(A_2/I_2)$-模生成 $I_2$。注意 $a_i\in I^{(1)}$ 且 $a_{in}\in I_n$。对每个 $n$，设 $J_n$ 是由 $a_{in}$（$i=1,\dots,k$）生成的 $A_n$ 的理想。由于 $\phi_{n2}\colon A_n\to A_2$（$n\ge 2$）是满射，有
\[
\phi_{n2}(I_n)=I_2=J_2=\phi_{n2}(J_n).
\]
由于 $\ker\phi_{n2}=I_n^2$，由此得
\[
I_n=J_n+I_n^2.
\]
结合 $I_n^n=0$，得
\[
I_n^{n-1}=J_n^{n-1}.
\]
设 $J^{(1)}$ 是由 $a_i$（$i=1,\dots,k$）生成的 $A$ 的理想。我们断言
\[
I^{(n)}=(J^{(1)})^n.
\]
于是 (i) 得证。我们有 $J^{(1)}\subset I^{(1)}$，从而 $(J^{(1)})^n\subset(I^{(1)})^n$。利用 $I_n^n=0$ 可证 $(I^{(1)})^n\subset I^{(n)}$。所以 $(J^{(1)})^n\subset I^{(n)}$。取
\[
x=(\dots,x_{n+2},x_{n+1},0,\dots,0)\in I^{(n)},
\]
则
\[
x_{n+1}\in\ker\phi_{n+1,n}=I_{n+1}^n=J_{n+1}^n.
\]
所以
\[
x_{n+1}=\sum_{i_1+\cdots+i_k=n}\lambda_{i_1\dots i_k,n+1}a_{1,n+1}^{i_1}\cdots a_{k,n+1}^{i_k},
\]
其中 $\lambda_{i_1\dots i_k,n+1}\in A_{n+1}$。选取 $\lambda_{i_1\dots i_k}\in A$，使它们在 $A_{n+1}$ 中的像为 $\lambda_{i_1\dots i_k,n+1}$。则
\[
x-\sum_{i_1+\cdots+i_k=n}\lambda_{i_1\dots i_k}a_1^{i_1}\cdots a_k^{i_k}\in I^{(n+1)}.
\]
用同样的论证，对每个 $i_1+\cdots+i_k=n+1$ 可找到 $\lambda_{i_1\dots i_k}\in A$，使
\[
x-\sum_{i_1+\cdots+i_k=n}\lambda_{i_1\dots i_k}a_1^{i_1}\cdots a_k^{i_k}
-\sum_{i_1+\cdots+i_k=n+1}\lambda_{i_1\dots i_k}a_1^{i_1}\cdots a_k^{i_k}\in I^{(n+2)}.
\]
类似地，对每个满足 $i_1+\cdots+i_k\ge n$ 的多重指标，可找到 $\lambda_{i_1\dots i_k}\in A$，使对每个 $m\ge 0$ 都有
\[
x-\sum_{n\le i_1+\cdots+i_k\le n+m}\lambda_{i_1\dots i_k}a_1^{i_1}\cdots a_k^{i_k}\in I^{(n+m+1)}.
\]
考虑级数
\[
\sum_{i_1+\cdots+i_k\ge n}\lambda_{i_1\dots i_k}a_1^{i_1}\cdots a_k^{i_k}.
\]
它是 $A=\plim_i A_i$ 中一个确定的元素，因为它在每个 $A_i$ 中的像都是有限和。我们有
\[
x-\sum_{i_1+\cdots+i_k\ge n}\lambda_{i_1\dots i_k}a_1^{i_1}\cdots a_k^{i_k}\in\bigcap_{m\ge 0}I^{(n+m+1)}=0.
\]
所以
\[
x=\sum_{i_1+\cdots+i_k\ge n}\lambda_{i_1\dots i_k}a_1^{i_1}\cdots a_k^{i_k}.
\]
可把 $\sum_{i_1+\cdots+i_k\ge n}\lambda_{i_1\dots i_k}a_1^{i_1}\cdots a_k^{i_k}$ 写成
\[
\sum_{i_1+\cdots+i_k=n}f_{i_1\dots i_k}(a_1,\dots,a_k)a_1^{i_1}\cdots a_k^{i_k},
\]
其中系数 $f_{i_1\dots i_k}(a_1,\dots,a_k)$ 是 $a_1,\dots,a_k$ 的形式幂级数，它们给出 $A$ 中确定的元素。于是
\[
x=\sum_{i_1+\cdots+i_k=n}f_{i_1\dots i_k}(a_1,\dots,a_k)a_1^{i_1}\cdots a_k^{i_k}\in(J^{(1)})^n.
\]
所以 $I^{(n)}\subset(J^{(1)})^n$，从而 $I^{(n)}=(J^{(1)})^n$。
\end{proof}

\begin{proposition}\label{it:1.5.11}
保留 \ref{it:1.5.10} 的假设与记号。设 $(M_n,\psi_{mn})_{n\in\bN}$ 是逆向系统，每个 $M_n$ 是 $A_n$-模，且对每对 $m\ge n$，$\psi_{mn}\colon M_m\to M_n$ 是与模结构及环同态 $\phi_{mn}\colon A_m\to A_n$ 相容的同态。设 $M_1$ 是有限生成 $A_1$-模，且每个 $\psi_{mn}$ 都是满射，其核为 $I_n^mM_m$。令 $M=\plim_n M_n$，并设 $\psi_n\colon M\to M_n$ 是典范同态。则 $M$ 是有限生成 $A$-模，且每个 $\psi_n$ 诱导同构
\[
M/(I^{(1)})^nM\cong M_n.
\]
\end{proposition}

\begin{proof}
设 $z_i=(\dots,z_{i2},z_{i1})$（$i=1,\dots,k$）是 $M=\plim_n M_n$ 中一族有限个元素，使得 $z_{i1}$（$i=1,\dots,k$）生成 $M_1$。对每个 $n$，设 $N_n$ 是由 $z_{in}$（$i=1,\dots,k$）生成的 $M_n$ 的子模。由于 $\psi_{n1}\colon M_n\to M_1$ 是满射，有
\[
\psi_{n1}(M_n)=M_1=N_1=\psi_{n1}(N_n).
\]
由于 $\ker\psi_{n1}=I_nM_n$，这蕴含
\[
M_n=N_n+I_nM_n.
\]
结合 $I_n^n=0$，得
\[
I_n^{n-1}M_n=I_n^{n-1}N_n.
\]
设 $N$ 是由 $z_i$（$i=1,\dots,k$）生成的 $M$ 的子模，并令 $M^{(n)}=\ker\psi_n$。我们断言
\[
M^{(n)}=I^{(n)}N.
\]
取 $n=0$，得 $M=N$，从而 $M$ 是有限生成 $A$-模。由于每个 $\psi_{mn}$ 都是满射，每个 $\psi_n$ 也是满射，且它诱导同构 $M/M^{(n)}\cong M_n$。但
\[
M/(I^{(1)})^nM=M/I^{(n)}M=M/M^{(n)}.
\]
命题得证。下面证明上述断言。对任意
\[
z=(\dots,z_{n+2},z_{n+1},0,\dots,0)\in M^{(n)},
\]
有
\[
z_{n+1}\in\ker\psi_{n+1,n}=I_{n+1}^nM_{n+1}=I_{n+1}^nN_{n+1}.
\]
所以
\[
z_{n+1}=\sum_{i=1}^{k}r_{i,n+1}z_{i,n+1},
\]
其中 $r_{i,n+1}\in I_{n+1}^n$。选取 $r_i^{(n+1)}\in A$，使它们在 $A_{n+1}$ 中的像为 $r_{i,n+1}$。则 $r_i^{(n+1)}\in I^{(n)}$，且
\[
z-\sum_{i=1}^{k}r_i^{(n+1)}z_i\in M^{(n+1)}.
\]
用同样的论证，可找到 $r_i^{(n+2)}\in I^{(n+1)}$，使
\[
z-\sum_{i=1}^{k}r_i^{(n+1)}z_i-\sum_{i=1}^{k}r_i^{(n+2)}z_i\in M^{(n+2)}.
\]
类似地，对每个 $m\ge 1$，可找到 $r_i^{(n+m)}\in I^{(n+m-1)}$（$i=1,\dots,k$），使
\[
z-\sum_{i=1}^{k}(r_i^{(n+1)}+r_i^{(n+2)}+\cdots+r_i^{(n+m)})z_i\in M^{(n+m)}.
\]
于是
\[
z=\sum_{i=1}^{k}\Bigl(\sum_{m=1}^{\infty}r_i^{(n+m)}\Bigr)z_i\in I^{(n)}N.
\]
所以 $M^{(n)}\subset I^{(n)}N$。但我们有 $M^{(n)}\supset I^{(n)}N$。因此 $M^{(n)}=I^{(n)}N$。
\end{proof}
\begin{tnote}
\ref{it:1.5.10} 与 \ref{it:1.5.11} 是纯粹代数的结论：它们说明，「$I$-进完全化」的逆向系统由它在第一层的有限生成模完全控制。这两条是本节的代数核心，也是 §2.7 中 Grothendieck 存在定理证明的出发点。[依据：原书 p.99–103 的命题陈述，以及本书 §2.7 对此处的引用]
\end{tnote}

设 $X$ 是诺特概形，$\cI$ 是 $\cO_X$ 的凝聚理想层。考虑环层空间 $(\widehat X,\cO_{\widehat X})$，其中 $\widehat X$ 是 $X$ 的闭子集
\[
\widehat X=\Supp(\cO_X/\cI),
\]
$\cO_{\widehat X}$ 是环层
\[
\cO_{\widehat X}=\plim_n(\cO_X/\cI^n)|_{\widehat X}.
\]
我们称 $(\widehat X,\cO_{\widehat X})$ 为 $X$ 关于 $\cI$ 的\emph{形式完全化}\idx{xingshiwanquanhua}{形式完全化}。常把 $(\widehat X,\cO_{\widehat X})$ 简记为 $\widehat X$。有环层空间的典范态射
\[
(i,i^\#)\colon(\widehat X,\cO_{\widehat X})\to(X,\cO_X),
\]
其中 $i\colon\widehat X\to X$ 是包含映射，且
\[
i^\#\colon\cO_X\to i_*\cO_{\widehat X}=\plim_n i_*i^{-1}(\cO_X/\cI^n)
\]
由典范态射族 $\cO_X\to i_*i^{-1}(\cO_X/\cI^n)$ 诱导。对每个 $n\in\bN$，设
\[
(X_n,\cO_{X_n})=\bigl(\Supp(\cO_X/\cI^n),(\cO_X/\cI^n)|_{\Supp(\cO_X/\cI^n)}\bigr)
\]
是由理想层 $\cI^n$ 定义的 $X$ 的闭子概形。则 $X_n$ 与 $\widehat X$ 有相同的底拓扑空间。设 $i_n\colon X_n\to X$ 与 $i_{nm}\colon X_n\to X_m$（$m\ge n$）是典范闭浸入。对每个 $n$，有典范态射
\[
\hat i_n\colon(X_n,\cO_{X_n})\to(\widehat X,\cO_{\widehat X}),
\]
其中 $\hat i_n\colon X_n\to\widehat X$ 是恒等映射，且
\[
\hat i_n^\#\colon\cO_{\widehat X}=\plim_n(\cO_X/\cI^n)|_{\widehat X}\to\cO_{X_n}=(\cO_X/\cI^n)|_{X_n}
\]
是典范态射。下图交换：
\begin{center}
\begin{tikzcd}
(X_n,\cO_{X_n}) \arrow[r,"i_n"] \arrow[d,"i_{nm}"'] \arrow[dr,"\hat i_n"] & (X,\cO_X)\\
(X_m,\cO_{X_m}) \arrow[r,"\hat i_m"'] & (\widehat X,\cO_{\widehat X}) \arrow[u,"i"']
\end{tikzcd}
\end{center}
\begin{tnote}
形式完全化只沿闭子集 $\widehat X$ 记录信息：它在 $\widehat X$ 上以 $\cI$ 的幂的逆向极限为结构层，在 $\widehat X$ 之外没有任何信息。这正是「沿闭子概形的形式邻域」的代数表达，也是 \ref{it:1.5.12} 要说明它是局部环层空间的原因。[依据：原书 p.103–104 的定义与 1.5.12]
\end{tnote}
